\section{Related work}
\label{sec:related}

\paragraph{Text-to-SQL benchmarks and the enterprise gap.}
Cross-domain text-to-SQL evaluation was standardized by
Spider~\cite{arxiv-1809-08887} and scaled by BIRD, which added dirty real-world
values, efficiency measurement, and the observation that external knowledge is
often required to bridge a question to a schema~\cite{arxiv-2305-03111}. Both
established the metric inherited here: execute predicted and gold SQL, compare
result sets. The gap to production analytics is documented from several
directions. Spider 2.0 moves to enterprise workflows spanning dialects and
multi-step transformation and reports that systems strong on earlier suites solve
a small fraction of it~\cite{arxiv-2411-07763}, with a multimodal sibling
covering the surrounding tooling~\cite{arxiv-2407-10956}. BEAVER builds questions
from an actual warehouse and finds models perform poorly even with prompt
engineering and retrieval, attributing the drop to complex schemas, complex
business questions, and private data that cannot have been trained
on~\cite{arxiv-2409-02038}. EntSQL isolates the knowledge dependency: across
1,066 aligned examples in five business domains, 96.0\% require information
beyond the question and schema, and the best evaluated system reaches 15.9\% with
long-form documents and 21.4\% with concise expert evidence, against a human
expert practical ceiling of 84.0\%~\cite{arxiv-2606-03363}. Recent suites push
further into business reasoning and decision
framing~\cite{arxiv-2510-07309,arxiv-2606-02109}, question whether SQL-string
accuracy is the right target~\cite{arxiv-2608-09254}, extend to database
operations agents~\cite{arxiv-2607-22165} and autonomous data-modeling
agents~\cite{arxiv-2606-19319}, and propose continuous benchmark generation to
keep pace with the systems~\cite{arxiv-2511-10049}.

The intervention studied here has a direct antecedent. Sequeda, Allemang, and
Jacob built an insurance-domain benchmark with an ontology and mappings defining
a knowledge graph over a SQL database, and reported zero-shot accuracy of 16\%
querying the SQL database directly against 54\% querying a knowledge-graph
representation of the same data~\cite{arxiv-2311-07509}. That result motivates
the question and also shows what a stronger design must add: a single contrast
between ``raw database'' and ``semantic representation'' moves access,
representation, and enforcement together, and the strength of the direct-SQL arm
bounds how much the comparison can establish. LiveSQLBench Large-v1 is used here because it
ships the intermediate rung: a structured, dependency-bearing knowledge base from
which both a searchable-raw-knowledge and a searchable-semantic-model condition
can be constructed from the same public source.

\paragraph{What execution accuracy does not measure.}
Pourreza and Rafiei show cross-domain text-to-SQL scores are sensitive to
evaluation choices, so reported gains can be artifacts of the comparison
procedure~\cite{arxiv-2310-18538}. SpotIt applies formal verification to the
evaluation itself and finds disagreements between result-set equality and query
equivalence~\cite{arxiv-2510-26840}, a question also attacked by asking whether
models can decide SQL equivalence~\cite{arxiv-2312-10321}. NL2SQL-BUGs targets
the complementary error, semantically wrong translations that survive
execution-based scoring~\cite{arxiv-2503-11984}; ROSE proposes an intent-centered
metric~\cite{arxiv-2604-12988}; SynSQL synthesizes databases to make
execution-based evaluation less sensitive to which rows happen to be
present~\cite{arxiv-2604-27261}. This is concrete here, because the official
LiveSQLBench Query scorer is lossy in ways we reproduced from its pinned source
(Section~\ref{sec:benchmark}). Rather than silently repair or silently inherit
the metric, we freeze both a faithful reimplementation and a corrected comparator
and commit in advance to reporting both.

\paragraph{Contamination and holdout reuse.}
The benchmark is public, so evaluated models may have seen it. Membership tests
for black-box models~\cite{arxiv-2310-17623}, pretraining-data
detection~\cite{arxiv-2310-16789}, leakage benchmarking~\cite{arxiv-2404-18824},
open contamination reporting~\cite{arxiv-2310-17589}, and analyses of what
training on a benchmark buys~\cite{arxiv-2409-01790} establish that contamination
is measurable and common; the practical mitigations are live or continuously
regenerated benchmarks~\cite{arxiv-2406-19314} and contamination-resistant
construction~\cite{arxiv-2505-08389}. We treat the dataset card's
contamination-free claim as an official claim, not an established result. The
mitigation available to us is different in kind: all four conditions draw on the
same public questions and comparisons are paired within question, so prior
exposure inflates all rungs together and the contrasts remain interpretable where
the absolute level does not. Development itself is adaptive data analysis. Dwork
et al. quantify the validity lost under holdout reuse and give mechanisms that
preserve it~\cite{arxiv-1506-02629}; Hardt analyzes ladder-style adaptive risk
estimation~\cite{arxiv-1706-02733}; Feldman, Frostig, and Hardt show how outcome
structure changes the overfitting rate~\cite{arxiv-1905-10360}; Chaibub Neto et
al. address the same problem on public
leaderboards~\cite{arxiv-1607-00091}. Hence the three-way partition of
Section~\ref{sec:design}.

\paragraph{Variance and the harness.}
Single-run benchmark numbers understate their own uncertainty. Miller sets out
error bars for language-model evaluations including clustered and paired
analyses~\cite{arxiv-2411-00640}; Madaan et al. quantify benchmark score variance
across seeds and configurations~\cite{arxiv-2406-10229}; Blackwell et al. measure
the reproducibility cost of nondeterminism~\cite{arxiv-2410-03492}; Coqueret et
al. state what randomness obliges researchers to report~\cite{arxiv-2607-24372};
Leech et al. catalogue the practices that inflate reported results, including
outcome selection after the fact~\cite{arxiv-2407-12220}. Separately, any claim
about a condition is a claim about a condition inside a harness. Zhang et al.
argue as a position that undisclosed harnesses make agent comparisons
uninterpretable~\cite{arxiv-2605-23950}; that paper reports no effect size, and
we attribute none to it. Two controlled studies bound the magnitude and disagree:
Vats and Golev measure a 0 to 8 percentage point paired pass-rate difference
across three coding harnesses but a 40-fold difference in tokens per solved
task~\cite{arxiv-2607-22585}, while Starace, preregistered, measures an accuracy
swing of up to 28 percentage points across three scaffolds on
GAIA~\cite{arxiv-2606-08529}; Gorinova et al. argue coding benchmarks are
misaligned with the agentic systems built on them~\cite{arxiv-2606-17799}. A 0 to
28 point range is too wide to assume away in either direction, which is why this
protocol neither assumes scaffold effects are negligible nor that they dominate,
and why C4 minus C3 is declared a system-level comparison whenever model parity
is unavailable.
